\chapter{Análisis Exploratorio de Datos (EDA)}

\section{Descripción General del Conjunto de Datos}
El conjunto de datos original cuenta con 48,895 registros y 16 variables que engloban tanto identificadores nominales, métricas temporales, como coordenadas geográficas continuas. Antes de someter estos datos a cualquier transformación predictiva, se realizó un Análisis Exploratorio de Datos (EDA) para diagnosticar problemas estructurales de asimetría, colinealidad y valores perdidos.

\section{Tratamiento de Valores Faltantes y Naturaleza MAR}
Una revisión de completitud reveló valores faltantes focalizados en cuatro columnas: \texttt{name}, \texttt{host\_name}, \texttt{last\_review} y \texttt{reviews\_per\_month}. 
Las dos primeras representan una pérdida negligente (menos del 0.05\%) y fueron imputadas con la etiqueta genérica "Unknown".

El caso de \texttt{last\_review} y \texttt{reviews\_per\_month} requirió un análisis diagnóstico más detallado. Se identificó que la falta de datos en estas variables no obedece a un error de recolección (MCAR, \textit{Missing Completely At Random}), sino que corresponde a una estructura MAR (\textit{Missing At Random}). Específicamente, existe una dependencia lógica determinista: estas variables son invariablemente nulas cuando la variable \texttt{number\_of\_reviews} es igual a 0. Esto significa que las propiedades que nunca han sido alquiladas o evaluadas carecen por definición de tasa de reseña mensual o fecha de última reseña. En base a esta evidencia, \texttt{reviews\_per\_month} fue imputado con 0, y \texttt{last\_review} fue transformada en una nueva variable derivada que cuantifica los días desde la última interacción, asignando un umbral máximo a los casos sin reseñas previas.

\section{Análisis de Distribuciones y Asimetría Severa}
El diagnóstico distribucional univariante alertó de desviaciones significativas respecto a la normalidad en la amplia mayoría de las variables predictoras continuas y en la variable objetivo (\texttt{price}). 

\begin{itemize}
    \item \textbf{Variable Objetivo (\texttt{price}):} La distribución del precio presenta un severo sesgo a la derecha (\textit{right skewness}), con una masa principal concentrada por debajo de los \$200 USD y valores atípicos extremos superando los \$10,000 USD. Para estabilizar la varianza y mejorar el comportamiento de los modelos lineales y paramétricos, se aplicó una transformación logarítmica $\log(1+x)$ (conocida como \texttt{log1p}).
    \item \textbf{Variables Temporales y de Actividad:} Las variables \texttt{minimum\_nights}, \texttt{number\_of\_reviews} y la derivada de recencia exhibieron un patrón de asimetría similar dictado por distribuciones de tipo Ley de Potencias. Por ende, la transformación \texttt{log1p} fue también aplicada a estos predictores.
\end{itemize}

Esta estandarización logarítmica mitiga sustancialmente la influencia desproporcionada de los valores atípicos sin necesidad de realizar recortes drásticos que suprimirían información valiosa (por ejemplo, listados genuinamente de lujo o penthouses).

\section{Distribución Espacial Inicial}
Visualmente, existe una hegemonía clara en los precios. El distrito de \textbf{Manhattan} ostenta los valores promedio más altos (aproximadamente \$196/noche), seguido de Brooklyn (\$124/noche), mientras que el Bronx y Staten Island fungen como áreas puramente residenciales y periféricas con valores significativamente menores. Asimismo, la variable \texttt{room\_type} polariza el mercado: las propiedades de tipo \textit{Entire home/apt} dominan casi con exclusividad la franja superior de precios.
